\documentclass[11pt]{article}

\usepackage[margin=1in]{geometry}
\usepackage[T1]{fontenc}
\usepackage[utf8]{inputenc}
\usepackage{lmodern}
\usepackage{microtype}
\usepackage{array}
\usepackage{booktabs}
\usepackage{enumitem}
\usepackage{tabularx}
\usepackage{xurl}
\usepackage[hidelinks]{hyperref}

\setlength{\parindent}{0pt}
\setlength{\parskip}{0.6em}
\setlength{\emergencystretch}{3em}
\setlist[itemize]{leftmargin=1.5em}
\setlist[enumerate]{leftmargin=1.75em}

\newcolumntype{Y}{>{\raggedright\arraybackslash}X}

\title{SEI Software Architecture Documentation\\[4pt]
\large A Practical Handout on the Views and Beyond Approach}
\author{}
\date{May 2026}

\begin{document}
\maketitle
\thispagestyle{empty}

\textbf{Purpose.} This handout explains the Software Engineering Institute (SEI) approach to software architecture documentation, commonly known as \emph{Views and Beyond} (V\&B). It also clarifies how an SEI-style software architecture document (SAD) should be understood in practice: as a structured documentation approach and reference model, not as a single universal standard template.\cite{sei-vb-factsheet,sei-vb-collection,sei-ref-standard}

\section*{What the SEI approach is}

SEI describes \emph{Views and Beyond} as a proven approach to documenting software architecture that uses the concept of a \emph{view} as the fundamental organizing principle for architecture documentation. A view represents a set of system elements and the relations among them.\cite{sei-vb-factsheet}

The goal is not to produce diagrams for their own sake. The goal is to capture an architecture in a form that stakeholders can use to build, analyze, test, operate, maintain, and evolve the system. SEI describes its documentation guidance as a way to produce a comprehensive documentation package that is useful to stakeholders and explicitly teaches the material in the context of other architecture description approaches, including ISO/IEC 42010.\cite{sei-book-page,sei-course}

\section*{What the SEI approach is not}

An SEI-style SAD should not be treated as:

\begin{itemize}
    \item a single mandatory industry standard for all software projects,
    \item a diagram dump with little explanation,
    \item a restatement of requirements or source code,
    \item or an exhaustive design encyclopedia that records every implementation detail.
\end{itemize}

Instead, the documentation should focus on architecturally significant decisions, the views needed by real stakeholders, the interfaces of architectural elements, and the rationale for major choices.\cite{sei-ref-standard,sei-interfaces}

\section*{Core ideas behind Views and Beyond}

\textbf{1. Choose views based on stakeholder concerns.} SEI's reference material requires explicit stakeholder representation and identification of stakeholder concerns. The document is therefore organized around the information different readers need, not around a one-size-fits-all chapter list.\cite{sei-ref-standard}

\textbf{2. Document the relevant views.} SEI's early documentation guidance states that documenting software architecture is primarily about documenting the relevant views and then augmenting them with trans-view information.\cite{sei-layers}

\textbf{3. Go beyond pictures.} Each view needs supporting explanation: element definitions, relations, constraints, behavior where relevant, interface information, and interpretation guidance. The view is the entry point, not the whole deliverable.\cite{sei-vb-factsheet,sei-interfaces}

\textbf{4. Document interfaces explicitly when they matter.} SEI's interface guidance states that interface documentation should record semantic as well as syntactic information. In other words, signatures alone are not enough; readers need to understand meaning, usage expectations, and reliance boundaries.\cite{sei-interfaces}

\textbf{5. Capture rationale.} A good SAD explains why significant architectural approaches were chosen, including why they were selected over other seriously considered options.\cite{sei-ref-standard}

\section*{What a practical SEI-style SAD usually contains}

SEI's 2005 reference standard for a software architecture document, written for DoD acquisition but still useful as a general reference structure, organizes the SAD around roadmap material, architectural background, views, relations among views, referenced materials, and directory information such as indexes and glossary content.\cite{sei-ref-standard}

\begin{itemize}
    \item \textbf{Documentation roadmap and overview:} helps readers find what they need quickly and states purpose, scope, document organization, and stakeholder representation.\cite{sei-ref-standard}
    \item \textbf{Architectural background:} records context, significant driving requirements, architectural approaches, analysis results, and requirements coverage.\cite{sei-ref-standard}
    \item \textbf{Views and view packets:} presents the architecture through one or more views, with each view packet documenting the elements, relations, interfaces, constraints, context, and local rationale for a bounded part of the system.\cite{sei-ref-standard}
    \item \textbf{Relations among views:} shows how elements in related views correspond and notes consistencies or known inconsistencies across views.\cite{sei-ref-standard}
    \item \textbf{Referenced materials:} provides complete citations to supporting material.\cite{sei-ref-standard}
    \item \textbf{Directory information:} includes an index, glossary, acronym list, and related reference aids.\cite{sei-ref-standard}
\end{itemize}

\section*{Views, viewpoints, and view packets}

A useful way to read the SEI approach is to separate three related concepts:

\begin{itemize}
    \item \textbf{Viewpoint:} the conventions for constructing and using a view.
    \item \textbf{View:} a representation of the system from the perspective of a related set of concerns.
    \item \textbf{View packet:} the smallest bundle of useful documentation that can be given to a stakeholder.\cite{sei-ref-standard}
\end{itemize}

This matters because a full architecture document should not be a monolith. It should be navigable. Different stakeholders should be able to go directly to the slices of architecture documentation that address their concerns.\cite{sei-ref-standard}

\section*{Interface documentation in the SEI approach}

SEI treats interface documentation as a first-class architectural concern. The interface report states that the guidance prescribes a standard organization for recording semantic as well as syntactic information about an interface. It also advises architects to include only information on which consumers are intended to rely.\cite{sei-interfaces}

In practice, that means an architect should document more than a method list, endpoint list, or message schema. A strong interface description should capture items such as:

\begin{itemize}
    \item what an element provides and requires,
    \item interaction assumptions and constraints,
    \item semantic meaning of operations or messages,
    \item usage expectations and error conditions,
    \item and any compatibility or reliance boundaries that downstream stakeholders must understand.\cite{sei-interfaces}
\end{itemize}

\section*{How to apply this on a real project}

A practical project-facing interpretation of the SEI approach looks like this:

\begin{enumerate}
    \item Identify stakeholder groups and the decisions or concerns each group needs the architecture documentation to support.\cite{sei-ref-standard}
    \item Select only the views that are necessary to address those concerns. Do not create views just because a template contains a heading for them.\cite{sei-layers,sei-vb-factsheet}
    \item For each selected view, define the element types, relations, constraints, and notation to be used.\cite{sei-ref-standard}
    \item Break large views into view packets that stakeholders can read and use without searching the entire document.\cite{sei-ref-standard}
    \item Document interfaces with both syntax and semantics.\cite{sei-interfaces}
    \item Record major architectural approaches, tradeoffs, and rationale, especially when a decision affects multiple views.\cite{sei-ref-standard}
    \item Maintain mappings among views so readers can trace how one architectural representation corresponds to another.\cite{sei-ref-standard}
    \item Keep the package current enough to remain operationally useful. A stale SAD quickly stops being architecture guidance and becomes archive material.\cite{sei-book-page}
\end{enumerate}

\section*{Recommended reading of the phrase ``SEI SAD''}

When practitioners refer to an ``SEI SAD,'' they are usually talking about one of two things:

\begin{itemize}
    \item an architecture document built using the \emph{Views and Beyond} method, or
    \item an SAD shaped by the SEI reference standard and related templates.
\end{itemize}

That shorthand is useful, but it should not be confused with a single universal SAD mandated by SEI for every context. The official SEI materials point instead to an approach, a collection of guidance, and downloadable templates and reference structures.\cite{sei-vb-collection,sei-template,sei-ref-standard}

\section*{Bottom line}

The SEI approach says that software architecture documentation should be organized around the right views for the right stakeholders, supported by interface documentation, rationale, and relations across views. The result should be a documentation package that helps people understand, build, evaluate, and sustain the system, rather than a document that merely looks complete.\cite{sei-vb-factsheet,sei-book-page,sei-interfaces}

\section*{References}

\begin{thebibliography}{9}

\bibitem{sei-vb-factsheet}
Software Engineering Institute.
\newblock \emph{Views and Beyond: The SEI Approach for Architecture Documentation}.
\newblock Fact sheet, February 14, 2018.
\newblock \url{https://www.sei.cmu.edu/library/views-and-beyond-the-sei-approach-for-architecture-documentation/}

\bibitem{sei-vb-collection}
Software Engineering Institute.
\newblock \emph{Views and Beyond Collection}.
\newblock SEI Library collection page.
\newblock \url{https://www.sei.cmu.edu/library/views-and-beyond-collection/}

\bibitem{sei-template}
Software Engineering Institute.
\newblock \emph{Views and Beyond Documentation Template}.
\newblock SEI Library page, February 6, 2006.
\newblock \url{https://www.sei.cmu.edu/library/views-and-beyond-documentation-template/}

\bibitem{sei-interfaces}
Felix Bachmann, Len Bass, Paul C. Clements, David Garlan, James Ivers, Reed Little, Robert Nord, and Judith Stafford.
\newblock \emph{Documenting Software Architecture: Documenting Interfaces}.
\newblock Technical Note CMU/SEI-2002-TN-015, Software Engineering Institute, June 2002.
\newblock \url{https://www.sei.cmu.edu/library/documenting-software-architecture-documenting-interfaces/}

\bibitem{sei-layers}
Felix Bachmann, Len Bass, Paul C. Clements, David Garlan, James Ivers, Reed Little, Robert Nord, and Judith Stafford.
\newblock \emph{Software Architecture Documentation in Practice: Documenting Architectural Layers}.
\newblock CMU/SEI-2000-SR-004, Software Engineering Institute, 2000.
\newblock \url{https://www.sei.cmu.edu/documents/5437/2000_003_001_13649.pdf}

\bibitem{sei-ref-standard}
John K. Bergey and Paul C. Clements.
\newblock \emph{Software Architecture in DoD Acquisition: A Reference Standard for a Software Architecture Document}.
\newblock Technical Note CMU/SEI-2005-TN-020, Software Engineering Institute, 2005.
\newblock \url{https://www.sei.cmu.edu/documents/2074/2005_004_001_14504.pdf}

\bibitem{sei-book-page}
Software Engineering Institute.
\newblock \emph{Documenting Software Architectures: Views and Beyond, 2nd Edition}.
\newblock SEI Library page, October 5, 2010.
\newblock \url{https://www.sei.cmu.edu/library/documenting-software-architectures-views-and-beyond-second-edition/}

\bibitem{sei-course}
Software Engineering Institute.
\newblock \emph{Documenting Software Architectures}.
\newblock SEI training course page.
\newblock \url{https://www.sei.cmu.edu/training/documenting-software-architectures/}

\end{thebibliography}

\end{document}



